\section{Persistence and Data Model}
\label{sec:schema}

\subsection{Storage Technology}

SQLite via GRDB is used rather than Core Data or SwiftData. The reasons are
specific and were weighed against the convenience cost:

\begin{itemize}
\item The schema is relational and query-heavy (library filtering, recipe
usage counts, cache eviction by LRU), which is SQL's home ground.
\item Migrations must be explicit and testable. The imaging model will evolve
across releases and a user's edit history must survive it; an inferred
migration is not acceptable for data the user cannot recreate.
\item Write patterns include high-frequency small writes (edit autosave) that
benefit from WAL mode and explicit transaction control.
\item The domain layer is already plain Swift values
(Section~\ref{sec:apparch}); an object graph framework would push framework
types into the domain, violating the layering of Table~\ref{tab:layers}.
\end{itemize}

The database is opened in WAL mode \cite{sqlitewal} with
\texttt{synchronous = NORMAL},
\texttt{foreign\_keys = ON}, and a 20\,MB page cache.

\subsection{Entity Relationship Overview}

\begin{figure}[htbp]
\centering
\begin{tikzpicture}[
  font=\scriptsize,
  ent/.style={draw=emLine, fill=emGray, rounded corners=2pt, minimum width=19mm, minimum height=6mm, align=center},
  rel/.style={draw=emLine!70, -{Latex[length=1.3mm]}},
  card/.style={font=\tiny, text=emLine, inner sep=1pt}
]
\node[ent] (roll)   at (0,3.2)   {roll};
\node[ent] (asset)  at (0,1.9)   {asset};
\node[ent] (edit)   at (0,0.6)   {edit};
\node[ent] (hist)   at (0,-0.7)  {edit\_history};
\node[ent] (rend)   at (3.3,1.9) {rendition};
\node[ent] (recipe) at (3.3,0.6) {recipe};
\node[ent] (stock)  at (3.3,-0.7){stock\_profile};
\node[ent] (coll)   at (3.3,3.2) {collection};
\node[ent] (memb)   at (1.65,3.2){collection\_item};

\draw[rel] (roll)  -- (asset)  node[card, midway, right] {1:N};
\draw[rel] (asset) -- (edit)   node[card, midway, right] {1:1};
\draw[rel] (edit)  -- (hist)   node[card, midway, right] {1:N};
\draw[rel] (asset) -- (rend)   node[card, midway, above] {1:N};
\draw[rel] (edit)  -- (recipe) node[card, midway, above] {N:1};
\draw[rel] (recipe)-- (stock)  node[card, midway, right] {N:1};
\draw[rel] (coll)  -- (memb)   node[card, midway, above] {1:N};
\draw[rel] (memb)  -- (asset.north) node[card, midway, left] {N:1};
\end{tikzpicture}
\caption{Entity relationships. \texttt{edit} carries the live recipe for an asset; \texttt{recipe} stores named, reusable parameter sets; \texttt{rendition} caches derived images keyed by content hash.}
\label{fig:erd}
\end{figure}

\subsection{Schema}

The complete DDL is given in Appendix~\ref{app:ddl}. The essential tables and
their design rationale follow.

\begin{lstlisting}[language=SQL, caption={Asset and roll tables.}, label={lst:ddl-asset}]
CREATE TABLE roll (
    id              TEXT PRIMARY KEY,          -- UUID
    name            TEXT NOT NULL,
    created_at      REAL NOT NULL,             -- Unix epoch, UTC
    -- A roll may carry a default recipe applied to new captures,
    -- which is how "shoot a roll of Portra" works as a workflow.
    default_recipe_id TEXT REFERENCES recipe(id) ON DELETE SET NULL,
    notes           TEXT
) STRICT;

CREATE TABLE asset (
    id              TEXT PRIMARY KEY,
    roll_id         TEXT REFERENCES roll(id) ON DELETE SET NULL,
    -- Exactly one of these two is non-null.
    file_relpath    TEXT,        -- managed store, relative to container
    ph_local_id     TEXT,        -- Photos library asset identifier
    content_hash    TEXT NOT NULL,             -- SHA-256 of source bytes
    original_uti    TEXT NOT NULL,
    pixel_width     INTEGER NOT NULL,
    pixel_height    INTEGER NOT NULL,
    captured_at     REAL,
    imported_at     REAL NOT NULL,
    -- Denormalized capture metadata for fast library filtering; the
    -- authoritative copy remains in the file.
    camera_model    TEXT,
    lens_model      TEXT,
    iso             INTEGER,
    exposure_time   REAL,
    f_number        REAL,
    focal_length    REAL,
    gps_lat         REAL,
    gps_lon         REAL,
    is_favorite     INTEGER NOT NULL DEFAULT 0,
    deleted_at      REAL,                      -- soft delete, 30-day purge
    CHECK ((file_relpath IS NULL) <> (ph_local_id IS NULL))
) STRICT;

CREATE INDEX idx_asset_roll     ON asset(roll_id, captured_at DESC);
CREATE INDEX idx_asset_captured ON asset(captured_at DESC) WHERE deleted_at IS NULL;
CREATE UNIQUE INDEX idx_asset_hash ON asset(content_hash) WHERE deleted_at IS NULL;
\end{lstlisting}

Three decisions in Listing~\ref{lst:ddl-asset} deserve comment. The
\texttt{CHECK} constraint enforcing exactly one storage location prevents the
ambiguous state where an asset is both managed and Photos-backed, which would
otherwise produce duplicate-detection bugs. The unique index on
\texttt{content\_hash} is partial, so soft-deleted assets do not block
re-import of the same file. And capture metadata is denormalized deliberately:
library filtering by camera or focal length must not require opening files.

\begin{lstlisting}[language=SQL, caption={Edit and history tables.}, label={lst:ddl-edit}]
CREATE TABLE edit (
    asset_id        TEXT PRIMARY KEY REFERENCES asset(id) ON DELETE CASCADE,
    -- The live recipe as canonical JSON. Stored as text rather than
    -- decomposed columns because the parameter set is versioned and
    -- evolves; a column-per-parameter schema would migrate on every
    -- imaging change.
    recipe_json     TEXT NOT NULL,
    recipe_version  INTEGER NOT NULL,          -- schema version of the JSON
    -- Hash of recipe_json; the rendition cache key.
    recipe_hash     TEXT NOT NULL,
    source_recipe_id TEXT REFERENCES recipe(id) ON DELETE SET NULL,
    updated_at      REAL NOT NULL
) STRICT;

CREATE TABLE edit_history (
    id              INTEGER PRIMARY KEY AUTOINCREMENT,
    asset_id        TEXT NOT NULL REFERENCES asset(id) ON DELETE CASCADE,
    seq             INTEGER NOT NULL,          -- position in the stack
    recipe_json     TEXT NOT NULL,
    changed_key     TEXT,                      -- ParameterKey raw value
    created_at      REAL NOT NULL
) STRICT;

CREATE UNIQUE INDEX idx_hist_asset_seq ON edit_history(asset_id, seq);
\end{lstlisting}

Storing the recipe as canonical JSON rather than as a wide table is the most
consequential schema decision in the application. The alternative --- a column
per parameter --- would require a migration on every imaging change, which will
happen many times. The JSON approach requires an in-application version field
and a migration chain for the JSON itself, which is code rather than schema, is
unit-testable, and can be applied lazily on read. The cost is that SQL cannot
query inside a recipe efficiently; this is mitigated by extracting the two
fields that actually need querying (\texttt{recipe\_hash} and
\texttt{source\_recipe\_id}) into real columns.

\begin{lstlisting}[language=SQL, caption={Recipe and stock profile tables.}, label={lst:ddl-recipe}]
CREATE TABLE recipe (
    id              TEXT PRIMARY KEY,
    name            TEXT NOT NULL,
    recipe_json     TEXT NOT NULL,
    recipe_version  INTEGER NOT NULL,
    is_builtin      INTEGER NOT NULL DEFAULT 0,
    -- Thumbnail rendered from a fixed reference image, so the recipe
    -- browser shows what each recipe does without an active asset.
    swatch_blob     BLOB,
    use_count       INTEGER NOT NULL DEFAULT 0,
    last_used_at    REAL,
    created_at      REAL NOT NULL,
    sort_order      INTEGER NOT NULL DEFAULT 0
) STRICT;

CREATE TABLE stock_profile (
    id              TEXT PRIMARY KEY,          -- e.g. "neg.portra400"
    kind            TEXT NOT NULL,             -- 'negative'|'print'|'mono'
    display_name    TEXT NOT NULL,
    -- Versioned so a profile correction can ship without silently
    -- changing the appearance of the user's existing edits; the edit
    -- records the profile version it was authored against.
    profile_version INTEGER NOT NULL,
    params_json     TEXT NOT NULL,             -- theta, A, C, grain, halation
    is_builtin      INTEGER NOT NULL DEFAULT 1,
    CHECK (kind IN ('negative','print','mono'))
) STRICT;
\end{lstlisting}

The \texttt{profile\_version} column addresses a real hazard. If a shipped
stock profile is later corrected --- and it will be, since profiles are fits to
measurement --- every existing edit made with that stock would silently change
appearance on update. \EM{} instead retains superseded profile versions and
binds each edit to the version it was authored against, offering an explicit
"update to corrected profile" action with a before/after preview. Silently
altering a user's finished images is a violation of the non-destructive promise
in FR-12, even though no file changed.

\begin{lstlisting}[language=SQL, caption={Rendition cache.}, label={lst:ddl-rend}]
CREATE TABLE rendition (
    id              INTEGER PRIMARY KEY AUTOINCREMENT,
    asset_id        TEXT NOT NULL REFERENCES asset(id) ON DELETE CASCADE,
    -- Cache key: (recipe_hash, size class, color space, profile versions)
    cache_key       TEXT NOT NULL,
    size_class      TEXT NOT NULL,   -- 'thumb'|'grid'|'screen'|'full'
    pixel_width     INTEGER NOT NULL,
    pixel_height    INTEGER NOT NULL,
    color_space     TEXT NOT NULL,
    file_relpath    TEXT NOT NULL,
    byte_size       INTEGER NOT NULL,
    created_at      REAL NOT NULL,
    last_access_at  REAL NOT NULL,
    CHECK (size_class IN ('thumb','grid','screen','full'))
) STRICT;

CREATE UNIQUE INDEX idx_rend_key ON rendition(cache_key);
CREATE INDEX idx_rend_lru ON rendition(last_access_at);
\end{lstlisting}

Cache eviction (NFR-12) is a straightforward LRU over
\texttt{idx\_rend\_lru}, with the constraint that \texttt{thumb} and
\texttt{grid} renditions for non-deleted assets are evicted last, since
regenerating them is what makes the library feel slow. Eviction runs on a
background task when total \texttt{byte\_size} exceeds 3.5\,GB, targeting
2.5\,GB.

\subsection{Migration Strategy}

Two independent version dimensions exist and must not be conflated:

\begin{enumerate}
\item \textbf{Schema version} --- the SQL structure, migrated by GRDB's
\texttt{DatabaseMigrator} with one registered, named, immutable migration per
change. Migrations are never edited after release; corrections ship as new
migrations.
\item \textbf{Recipe JSON version} --- the parameter model, migrated in Swift
by a chain of pure functions $f_n: \mathrm{JSON}_n \to \mathrm{JSON}_{n+1}$.
\end{enumerate}

Recipe migration is lazy: a recipe is migrated on read and written back only
when the asset is next edited. This avoids a long blocking migration on update
for users with large libraries, which is a common source of one-star reviews.

The migration chain is tested by a golden-file suite: for each historical
version, a corpus of recipes is stored, migrated forward to current, rendered,
and compared against a stored reference image with a $\Delta E_{00}$
\cite{sharma2005} tolerance.
This is the only mechanism that reliably catches a migration that parses
successfully but changes appearance.

\begin{lstlisting}[language=Swift, caption={Migration registration.}]
var migrator = DatabaseMigrator()
#if DEBUG
migrator.eraseDatabaseOnSchemaChange = true    // development only
#endif

migrator.registerMigration("v1_initial") { db in
    try db.execute(sql: ddlV1)
}
migrator.registerMigration("v2_rendition_lru") { db in
    try db.alter(table: "rendition") { t in
        t.add(column: "last_access_at", .double).notNull().defaults(to: 0)
    }
    try db.create(index: "idx_rend_lru", on: "rendition", columns: ["last_access_at"])
}
migrator.registerMigration("v3_profile_versioning") { db in
    try db.alter(table: "stock_profile") { t in
        t.add(column: "profile_version", .integer).notNull().defaults(to: 1)
    }
}
\end{lstlisting}

\subsection{Sidecar Format}

FR-17 requires an interchangeable recipe file. The format is JSON with a
documented schema:

\begin{lstlisting}[language=Swift, caption={Sidecar structure (illustrative).}]
{
  "emulsion":       { "format": 1, "app": "1.0.0" },
  "recipeVersion":  4,
  "name":           "Portra 400 / 2383, warm",
  "negativeStock":  { "id": "neg.portra400", "version": 3 },
  "printStock":     { "id": "prt.2383",      "version": 2 },
  "chemistry":      { "id": "chem.c41",      "version": 1 },
  "format":         "135",
  "seed":           874112035,
  "parameters": {
    "pushPull":        0.0,
    "printerLightR":  25, "printerLightG": 25, "printerLightB": 27,
    "printDensity":    3,
    "grainAmount":     1.0,
    "halationIntensity": 0.22
  },
  "provenance": {
    "createdAt": "2026-03-14T09:22:41Z",
    "sourceAsset": "5C1B...",
    "note": "matched to lab scan of roll 04 frame 12"
  }
}
\end{lstlisting}

Only parameters that differ from the resolved defaults are written, which keeps
sidecars small and, more importantly, means a sidecar authored against one stock
applies sensibly to another --- the unspecified parameters take the new stock's
defaults rather than the old stock's values. This is the correct semantics for a
"look" and it falls directly out of the sparse-recipe design of
Section~\ref{sec:apparch}.
